\VerdierPhysicalPageAnchor{0081}%
\VerdierSemanticAnchor{ch1-2-5-11}
\noindent\textbf{2.5.11.}\quad
\VerdierSemanticAnchor{ch1-2-5-11-homotopism}
设 $X$ 和 $Y$ 是 $C$ 的两个复形。若 $\operatorname{comp}(C)$ 中的态射
$f\colon X\longrightarrow Y$ 在 $K(C)$ 中的像 $\dot f$ 是同构，则称
$f$ 为\emph{同伦等价}。
\VerdierSemanticAnchor{ch1-2-5-11-homotopic-complexes}
若复形 $X$ 和 $Y$ 在 $K(C)$ 中同构，则称它们\emph{同伦等价}。

\VerdierSemanticAnchor{ch1-2-5-12}
\noindent\textbf{2.5.12.}\quad
\VerdierSemanticAnchor{ch1-2-5-12-setup}
设 $X$ 和 $Y$ 是 $C$ 的两个复形，
$f,g\colon X\rightrightarrows Y$ 是 $\operatorname{comp}(C)$ 中的两个态射，
$a,b\colon f\rightsquigarrow g$ 是连接 $f$ 与 $g$ 的两个同伦。
\VerdierSemanticAnchor{ch1-2-5-12-second-order-homotopy}
连接 $a$ 与 $b$ 的一个\emph{二阶同伦}，记作
$\alpha\colon a\rightsquigarrow b$，是一个态射族
\[
  \alpha=(\alpha^i\colon X^{i+2}\longrightarrow Y^i,\ i\in\mathbb Z)
\]
并满足：

\VerdierSemanticAnchor{ch1-2-5-12-eq1}
\begin{equation}
  b^i-a^i
  =\alpha^i d_X^{i+1}-d_Y^{i-1}\alpha^{i-1}.
  \tag{2.5.12.1}
\end{equation}

\VerdierSemanticAnchor{ch1-2-5-13}
\noindent\textbf{2.5.13.}\quad
最后指出，同构 (2.4.1)

\VerdierSemanticAnchor{ch1-2-5-13-complex-opposite-isomorphism}
\begin{equation*}
  \iota_1\colon
  \operatorname{comp}(C)^\circ
  \longrightarrow
  \operatorname{comp}(C^\circ)
\end{equation*}
由 (2.5.7) 唯一确定一个同构

\VerdierSemanticAnchor{ch1-2-5-13-homotopy-opposite-isomorphism}
\begin{equation*}
  \iota_1\colon
  K(C)^\circ
  \longrightarrow
  K(C^\circ),
\end{equation*}
\VerdierSemanticAnchor{ch1-2-5-13-z-category-isomorphism}
它是 $\mathbb Z$-范畴的同构。
\VerdierSemanticAnchor{ch1-2-5-13-identification}
今后我们通常借此同构将范畴 $K(C)^\circ$ 与 $K(C^\circ)$ 视为同一范畴。
